\documentclass[12pt,a4paper]{article}

\usepackage{cmap}
\usepackage[T2A]{fontenc}
\usepackage[utf8]{inputenc}
\usepackage[russian]{babel}
\usepackage{amsmath,amssymb}
\usepackage[a4paper,margin=2cm]{geometry}

\setlength{\parindent}{0pt}
\setlength{\parskip}{0.45em}
\sloppy

\begin{document}

\begin{center}
{\Large\bfseries Билет 5: вопросы для проверки знаний}
\end{center}

\noindent\fbox{%
  \begin{minipage}{\dimexpr\linewidth-2\fboxsep-2\fboxrule\relax}
  \normalfont
Определение измеримости по Жордану множества в $n$-мерном евклидовом
пространстве. Критерий измеримости. Измеримость объединения, пересечения и
разности измеримых множеств. Конечная аддитивность меры Жордана.
  \end{minipage}%
}

\section*{Вопросы на понимание}

\begin{enumerate}
\item Чем нижняя мера Жордана отличается от верхней: какое приближение множества использует каждая из них?
\item Что геометрически означает условие, что для любого $\varepsilon>0$ можно найти клеточные множества $A\subset E\subset B$ с $\mu(B)-\mu(A)<\varepsilon$?
\item Почему в определении измеримости по Жордану требуется конечность общей меры, а не только равенство нижней и верхней мер?
\item Почему в критерии измеримости по границе условие ограниченности нельзя опустить? Рассмотрите пример $\mathbb{R}^n$.
\item Почему множества $(0,1)$, $[0,1]$ и $(0,1]$ имеют одну и ту же меру Жордана?
\item Приведите пример ограниченного множества в $\mathbb{R}$, которое не измеримо по Жордану, и объясните это через его границу.
\item Почему граница меры нуль означает, что внутреннее и внешнее клеточные приближения могут отличаться сколь угодно малой мерой?
\item Почему из измеримости $E$ и $F$ следует измеримость $E\cup F$, $E\cap F$ и $E\setminus F$? Какую роль играют включения для границ?
\item При образовании разности $E\setminus F$ могут появиться новые граничные точки. Почему они все равно контролируются границами $E$ и $F$?
\item Почему в конечной аддитивности достаточно требовать, чтобы множества не имели общих внутренних точек, а не обязательно были непересекающимися?
\item Что нарушится в формуле $\mu(E\cup F)=\mu(E)+\mu(F)$, если $E$ и $F$ пересекаются по множеству с ненулевой мерой?
\item Почему конечная аддитивность согласуется с тем, что клетку можно разбить на конечное число меньших клеток без изменения общей меры?
\end{enumerate}

\section*{Источники для сверки}

\texttt{target/answer\_question\_05.tex};
\texttt{IvGE\_1.pdf}, стр. 200--206.

\end{document}
